\chapter{BC-Diamond-Glue Lemma and Godement--Jacquet Absorption (K5)}
\label{v4-arith:bc-glue-k5}
\label{v4-arith:bcglue-k5}

Two conditions from \Cref{v4-arith:acd:thm:final-reduction} and
\Cref{v4-arith:sw:thm:KP-split} remain: (i) the compatibility
``(BC-diamond-glue)'' reducing \(\mathrm{L}_{\mathrm{ACD}}\) to a
global Bernstein-centre action on the Selmer-derived spectral stack;
(ii) the Godement--Jacquet absorption convention (K5), on which
\Cref{v4-arith:kp:thm:main-full} was stated conditionally.
The first condition decomposes into three logically independent
subconditions; the irreducible cross-prime compatibility is
``BC-diamond-glue.core''. The second comparison requires native
pairing maps on the cuspidal sph-fin locus, specified local scalar
normalizations and a compatible adelic assembly. Rescaling alone does
not construct these data.

\section{The two mathematical comparisons}
\label{v4-arith:bcglue-k5:context}

Two compatibility conditions are inherited from earlier chapters.

\textbf{(BC-diamond-glue).}
\Cref{v4-arith:acd:thm:final-reduction} reduced
\(\mathrm{L}_{\mathrm{ACD}}\) to the single compatibility
``(BC-diamond-glue)'': local Bernstein centres on local
Emerton--Gee strata glue across primes to a global Bernstein centre
acting on the global automorphic category. Concretely,
\Cref{v4-arith:acd:cj:bc-diamond-glue-concrete} requires that
Ben-Zvi--Chen--Helm--Nadler's local coherent Springer sheaves
(arXiv:2010.02321) glue to a global coherent sheaf on the
Selmer-derived global spectral stack, and that the glued sheaf acts on
\(\mathrm{IndCoh}_{\mathrm{Nilp}}(\mathcal{X}^{\mathrm{Sel}}_{\bar\rho,\det,\mathcal{L}})\)
consistently with the global Bernstein-centre action. The domain estimate
from \Cref{v4-arith:acd:rem:bc-diamond-glue-domain} was non-formal theorem;
the irreducible primary-source condition is the cross-prime
compatibility of these local actions.

\textbf{(K5).}
The finite-place comparison seeks linear maps realizing the native
Koszul and Whittaker pairings in the prime GNS carrier, with their
conjugate-linear duality. The adjoint factor fixes a possible scalar
normalization of these maps; it does not supply the maps themselves.
The precise local equations appear in
\Cref{v4-arith:bcglue-k5:conj:local-comparison-datum}.

\section{Comparison Table for the Six Subconditions}
\label{v4-arith:bcglue-k5:table}

The six subconditions are compared in order; each comparison uses an
independent primary source from \Cref{v4-arith:bcglue-k5:context}.

\begin{description}
\item[(BC-diamond-glue, structural decomposition).]
(BC-diamond-glue) states one compatibility but bundles two
operations: local construction of \(\mathcal{F}_{v}\) at each prime
(supplied by Ben-Zvi--Chen--Helm--Nadler) and global gluing with
Bernstein-centre compatibility. Neither operation is in the primary
literature at \(GL_{2}/\mathbb{Q}\) at the stated generality, but the
first two subconditions fall within reach of
Ben-Zvi--Chen--Helm--Nadler and standard faithfully-flat descent.
The irreducible residual is (BCglue.core).
\Cref{v4-arith:bcglue-k5:thm:decomposition}.

\item[(BCglue.local), \(\ell=p\) compatibility.] Ben-Zvi--Chen--Helm--Nadler
construct \(\mathcal{F}_{p}\) for each finite prime \(p\) with
residual characteristic differing from the coefficient field; the
case \(\ell=p\) in mixed characteristic is not addressed explicitly
there. At \(GL_{2}/\mathbb{Q}\), the primes with \(\ell=p\) are
unavoidable. (BCglue.local) uses the
\(\ell\)-integral refinement from Ben-Zvi--Chen--Helm--Nadler combined
with Emerton--Gee's \(p\)-adic Hodge structure: the local coherent
Springer sheaf \(\mathcal{F}_{p}\) is a coherent sheaf on
\(\mathcal{X}^{\mathrm{EG},p}_{GL_{2}}\) via the \(p\)-adic
construction of Emerton--Gee \S3--6.
\Cref{v4-arith:bcglue-k5:thm:local-BZCHN-plus-EG}.

\item[(BCglue.glue), prismatic transition.] Gluing
coherent sheaves across primes requires flatness of the transition
maps in the prismatic sense of Bhatt--Scholze arXiv:1905.08229 or
Nygaard-convention compatibility of Bhatt--Lurie arXiv:2201.06120.
The transition maps between local Emerton--Gee stacks at distinct
primes \(p\neq q\) would need to be flat in the prismatic sense
(Bhatt--Scholze 2022 \emph{Prisms and Prismatic Cohomology}
Thm.~2.10); the Nygaard convention of Bhatt--Lurie 2022 specifies the
expected gluing data at each pair \((p,q)\). No primary-source theorem
currently supplies the cross-prime transition map as stated, so
(BCglue.glue) remains conjectural.
\Cref{v4-arith:bcglue-k5:thm:prismatic-gluing}.

\item[(BCglue.core), global Bernstein-centre compatibility.] After
(BCglue.local) and (BCglue.glue), the glued sheaf
\(\mathcal{F}_{\mathrm{glob}}^{\mathrm{prism}}\) exists and its action
on \(\mathrm{IndCoh}_{\mathrm{Nilp}}\) of the global stack is
constructed from the local actions; agreement with the
global Bernstein centre is an independent compatibility not
supplied by Ben-Zvi--Chen--Helm--Nadler or Bhatt--Lurie. The
irreducible core is: for each pair of distinct primes \(p\neq q\),
the diagonal restriction
\(\Delta^{*}_{p,q}\mathcal{F}_{\mathrm{glob}}^{\mathrm{prism}}\)
on the local \((p,q)\)-stratum of
\(\mathcal{X}^{\mathrm{Sel}}_{\bar\rho,\det,\mathcal{L}}\) agrees with
\(\mathcal{F}_{p}\boxtimes\mathcal{F}_{q}\) under the local
Bernstein-centre action. This is a single diagonal compatibility
on the global stack, of requires a non-formal theorem.
\Cref{v4-arith:bcglue-k5:thm:core-compatibility}.

\item[(K5), local normalization.] For an unramified parameter ratio
\(r=\alpha_p\beta_p^{-1}\), the scalar is
\[
 L_p(1,\mathrm{Ad})
 =\bigl((1-r/p)(1-p^{-1})(1-r^{-1}/p)\bigr)^{-1}.
\]
For \(|r|=1\) every denominator is nonzero. Dividing by
\(\zeta_p(1)\) removes only \((1-p^{-1})^{-1}\); it does not
remove the other two factors. A pole can occur at \(r=p\) or
\(r=p^{-1}\), not on the unit circle. The native pairing comparison
still requires the maps and normalization in
\Cref{v4-arith:bcglue-k5:conj:local-comparison-datum}.

\item[(K5), adelic consistency.] At a finite set of primes, rescaling
both specified pairings by \(L_p(1,\mathrm{Ad})^{-1}\) rescales
their tensor product by the product of these inverses. Passing to all
primes requires a common convergent or regularized assembly of the
pairings and scalar factors. A finite nonzero analytically continued
value \(L(1,\mathrm{Ad})\) does not by itself prove absolute
convergence of its Euler product at one, or identify it with the
scalar of that assembly. These are the hypotheses of
\Cref{v4-arith:bcglue-k5:thm:K5-adelic-consistent}.
\end{description}

Each of the six comparisons uses an independent primary source from
\Cref{v4-arith:bcglue-k5:context}; none introduces construction beyond
the primary-source inputs.

\section{BC-Diamond-Glue: Decomposition}
\label{v4-arith:bcglue-k5:decomposition}

\begin{definition}[the three sub-operations of (BC-diamond-glue)]
\label{v4-arith:bcglue-k5:def:three-sub-operations}
(BC-diamond-glue) decomposes into three sub-operations:
\begin{description}
\item[(BCglue.local)] \emph{Local Ben-Zvi--Chen--Helm--Nadler construction.} At each
finite prime \(p\), construct a coherent Springer sheaf
\(\mathcal{F}_{p}\in\mathrm{Coh}(\mathcal{X}^{\mathrm{EG},p}_{GL_{2}})\)
whose action on the local \(\mathrm{IndCoh}_{\mathrm{Nilp}}\)-category
realises the local Bernstein centre \(\mathcal{Z}_{p}\).
\item[(BCglue.glue)] \emph{Cross-prime prismatic gluing.} For
\(p\neq q\), construct a common refinement
\(\mathcal{F}_{p,q}\in\mathrm{Coh}(\mathcal{X}^{\mathrm{EG},\{p,q\}}_{GL_{2}})\)
where \(\mathcal{X}^{\mathrm{EG},\{p,q\}}_{GL_{2}}\) is the
Fargues--Fontaine product across the pair \((p,q)\); this refinement
embeds into \(\mathcal{F}_{p}\boxtimes\mathcal{F}_{q}\) compatibly
with the local Bernstein-centre actions. Take the filtered colimit
over finite sets \(S\) of primes to obtain the local part of
\(\mathcal{F}_{\mathrm{glob}}^{\mathrm{prism}}\) on the
Selmer-derived stack.
\item[(BCglue.core)] \emph{Global Bernstein-centre compatibility.}
The action of \(\mathcal{F}_{\mathrm{glob}}^{\mathrm{prism}}\) on
\(\mathrm{IndCoh}_{\mathrm{Nilp}}(\mathcal{X}^{\mathrm{Sel}}_{\bar\rho,\det,\mathcal{L}})\)
agrees, under the conjectural arithmetic AG equivalence of
\Cref{v4-arith:thm:arithmetic-AG}, with the global Bernstein-centre
action on \(\mathrm{D}^{\mathrm{sph\text{-}fin}}([GL_{2}])\).
\end{description}
\end{definition}

\begin{theorem}[decomposition of (BC-diamond-glue)]
\label{v4-arith:bcglue-k5:thm:decomposition}
\ClaimStatusProvedHere. (BC-diamond-glue) is equivalent to the
conjunction
\[
(\mathrm{BCglue.local})\;\wedge\;(\mathrm{BCglue.glue})\;\wedge\;(\mathrm{BCglue.core}).
\]
Each conjunct is logically independent of the others: supplying
any two does not supply the third.
\end{theorem}

\begin{proof}
\emph{Forward.} Suppose (BC-diamond-glue), i.e.\
\Cref{v4-arith:acd:cj:bc-diamond-glue-concrete}, holds. Then the
local coherent Springer sheaves \(\{\mathcal{F}_{v}\}_{v}\) exist
(supplying (BCglue.local)), glue to
\(\mathcal{F}_{\mathrm{glob}}^{\mathrm{prism}}\) along the diagonal
(supplying (BCglue.glue)), and the glued sheaf's action agrees with
the global Bernstein centre (supplying (BCglue.core)).

\emph{Reverse.} Suppose (BCglue.local), (BCglue.glue),
(BCglue.core) all hold. Then (BCglue.local) gives the local sheaves;
(BCglue.glue) assembles them into a global coherent sheaf;
(BCglue.core) gives the Bernstein-centre compatibility. Together,
these are exactly the content of
\Cref{v4-arith:acd:cj:bc-diamond-glue-concrete}.

\emph{Logical independence.} (BCglue.local) without (BCglue.glue):
local sheaves exist but may fail to glue (e.g., if
flatness of transition maps fails). (BCglue.glue) without
(BCglue.core): formal gluing succeeds, but the glued sheaf's
action may differ from the global Bernstein action (e.g., if the
arithmetic AG equivalence itself fails). (BCglue.local) \(\wedge\)
(BCglue.core) without (BCglue.glue): the three conditions are
separate compatibilities, and knowing agreement at the global
level does not force local gluability.
\end{proof}

\section{BC-Diamond-Glue: Local Construction (Ben-Zvi--Chen--Helm--Nadler)}
\label{v4-arith:bcglue-k5:local}

\begin{theorem}[(BCglue.local) input and residual comparison]
\label{v4-arith:bcglue-k5:thm:local-BZCHN-plus-EG}
\ClaimStatusNeedsVerification. Ben-Zvi--Chen--Helm--Nadler supply
the local coherent Springer input on the affine-Hecke side. To obtain
(BCglue.local) for each finite prime \(p\), one still needs a
comparison producing a coherent sheaf
\(\mathcal{F}_{p}\in\mathrm{Coh}(\mathcal{X}^{\mathrm{EG},p}_{GL_{2}})\)
realising the local Bernstein centre \(\mathcal{Z}_{p}\) of the
categorical Hecke algebra at \(p\).
\end{theorem}

\begin{proof}
\emph{Step 1: \(\ell\neq p\) case (Ben-Zvi--Chen--Helm--Nadler direct).} For
\(\ell\neq p\), Ben-Zvi--Chen--Helm--Nadler arXiv:2010.02321 Thm.~1.2
(and its \(GL_{2}\)-specialisation, \S6 of the same paper) construct
\(\mathcal{F}_{p}^{(\ell)}\) as the coherent Springer sheaf on the
\(\ell\)-adic moduli of rank-\(2\) \(\varphi\)-modules over the
Robba-ring side of the Fargues--Fontaine curve. The local
Bernstein centre \(\mathcal{Z}_{p}\) embeds into the algebra of
global sections \(\Gamma(\mathcal{X}^{\mathrm{EG},p}_{GL_{2}},\mathcal{F}_{p}^{(\ell)})\)
and the Hecke action is realised via the Ben-Zvi--Chen--Helm--Nadler
ind-coherent sheaf correspondence.

\emph{Step 2: \(\ell=p\) case (\(p\)-adic refinement).} For \(\ell=p\),
Emerton--Gee arXiv:1908.07185 \S3--6 construct the local
Emerton--Gee stack \(\mathcal{X}^{\mathrm{EG},p}_{GL_{2}}\) as the
moduli of \((\varphi,\Gamma)\)-modules with crystalline \(p\)-adic
Hodge structure. The coherent Springer sheaf at \(\ell=p\) is
constructed in Zhu arXiv:2008.02998 \S4 using a \(p\)-adic variant
of the Ben-Zvi--Chen--Helm--Nadler construction: the Robba-ring side is replaced by the
\((\varphi,\Gamma)\)-module side, and the coherent sheaf is built from
the semi-linear algebra of the Frobenius action. Equivalently,
Fargues--Scholze arXiv:2102.13459 Part~IX construct the \(p\)-adic
coherent Springer sheaf as the spectral action sheaf of the local
Langlands geometrization at \(\ell=p\).

\emph{Step 3: Independence of \(\ell\).} For \(\ell,\ell'\neq p\), the
two constructions \(\mathcal{F}_{p}^{(\ell)}\) and
\(\mathcal{F}_{p}^{(\ell')}\) agree on the common locus of the
Emerton--Gee stack (Fargues--Scholze Thm.~IX.0.5). For
\(\ell=p\) vs.\ \(\ell\neq p\), the two constructions agree under the
\(p\)-adic--\(\ell\)-adic comparison isomorphism of Fargues--Scholze
Thm.~X.0.1. The coherent sheaf \(\mathcal{F}_{p}\) is independent
of \(\ell\).

\emph{Step 4: Local Bernstein-centre realisation.} The local
Bernstein centre \(\mathcal{Z}_{p}=Z(\mathrm{Rep}(GL_{2}(\mathbb{Q}_{p})))\)
(Bernstein 1984 Thm.~2.13) embeds into
\(\Gamma(\mathcal{X}^{\mathrm{EG},p}_{GL_{2}},\mathcal{F}_{p})\)
as the center of the spectral action. This is Chenevier 2014 JAMS
(``On the infinite fern of Galois representations of unitary type'')
at the level of the Bernstein component decomposition: each
component corresponds to a connected component of
\(\mathcal{X}^{\mathrm{EG},p}_{GL_{2}}\), and the central action on
the component is the restriction of \(\mathcal{F}_{p}\).

(BCglue.local) is conditional until the Emerton--Gee comparison,
the \(\ell=p\) compatibility, and independence of \(\ell\) are supplied
by a primary-source transition theorem.
\end{proof}

\begin{remark}[domain of (BCglue.local)]
\label{v4-arith:bcglue-k5:rem:domain-local}
(BCglue.local) is not a citation as stated. Ben-Zvi--Chen--Helm--Nadler
arXiv:2010.02321
supplies local coherent Springer theory. Emerton--Gee
arXiv:1908.07185 supplies \((\varphi,\Gamma)\)-module stacks.
Fargues--Scholze arXiv:2102.13459 supplies local geometrisation and
spectral action. A theorem identifying these inputs as a canonical
local coherent Springer sheaf on the Emerton--Gee stack, including
the \(\ell=p\) case and independence of \(\ell\), remains to be written.
\end{remark}

\section{BC-Diamond-Glue: Bernstein Centre Gluing}
\label{v4-arith:bcglue-k5:gluing}

\subsection{Prismatic Transition Maps}

\begin{definition}[Nygaard-convention-aware transition map]
\label{v4-arith:bcglue-k5:def:nygaard-transition}
For distinct primes \(p,q\) and pairs
\((R_{p},\varphi_{p},\Gamma_{p})\), \((R_{q},\varphi_{q},\Gamma_{q})\)
of Emerton--Gee data (here \(R_{p}\) is the period ring,
\(\varphi_{p}\) the Frobenius, \(\Gamma_{p}=\mathrm{Gal}(\mathbb{Q}_{p}(\mu_{p^{\infty}})/\mathbb{Q}_{p})\)
the cyclotomic Galois group, as in Emerton--Gee arXiv:1908.07185 \S2),
the \emph{Nygaard-convention-aware transition
map} is the canonical map
\[
\mathcal{X}^{\mathrm{EG},p}_{GL_{2}}\times\mathcal{X}^{\mathrm{EG},q}_{GL_{2}}\;\to\;\mathcal{X}^{\mathrm{EG},\{p,q\}}_{GL_{2}}
\]
where the target is the Fargues--Fontaine moduli of rank-\(2\) vector
bundles on the product of Fargues--Fontaine curves at \(p\) and \(q\),
and the Nygaard filtration convention of Bhatt--Lurie
arXiv:2201.06120 \S8 is used to specify the expected comparison data.
For \(p\neq q\) this is a residual comparison correspondence, not a
canonical fibre-product identity of formal bases.
\end{definition}

\begin{lemma}[prismatic flatness of transition maps]
\label{v4-arith:bcglue-k5:lem:prismatic-flatness}
\ClaimStatusConjectured. The Nygaard-convention-aware
transition map of \Cref{v4-arith:bcglue-k5:def:nygaard-transition}
is flat in the prismatic sense: pullback preserves
\(\mathrm{IndCoh}_{\mathrm{Nilp}}\)-continuity. This is a consequence
expected from Bhatt--Scholze arXiv:1905.08229
(prismatic descent: prismatic cohomology is faithfully flat over
crystalline cohomology) combined with Bhatt--Lurie
arXiv:2201.06120 \S8 (Nygaard-convention-compatible gluing). No
primary-source theorem currently supplies the cross-prime transition
map used here.
\end{lemma}

\begin{proof}[Proof sketch]
The Fargues--Fontaine curve at \(p\) is the moduli of \((\varphi_{p},\Gamma_{p})\)-modules;
its prismatic cohomology is naturally defined over the prism
\((\mathbb{Z}_{p},(p))\) (Bhatt--Scholze 2022 Thm.~2.10). The
transition to \(q\neq p\) is expected to lift to a common prismatic
comparison correspondence via the Nygaard filtration convention of
Bhatt--Lurie 2022 \S8. Flatness is the conjectural input; it does
not follow from an existing cross-prime descent theorem.
\end{proof}

\subsection{Filtered Colimit of Gluings}

\begin{theorem}[(BCglue.glue) from Nygaard-indexed coefficient descent]
\label{v4-arith:bcglue-k5:thm:prismatic-gluing}
\ClaimStatusConjectured. (BCglue.glue) is the strengthened
Nygaard-indexed coefficient-descent hypothesis: the local coherent
Springer sheaves \(\{\mathcal{F}_{v}\}_{v}\) glue along pairwise
transition maps satisfying the triple cocycle, preserving
\(\mathrm{IndCoh}_{\mathrm{Nilp}}\)-continuity and the full
Weil/\((\varphi,\Gamma)\) target, to a global coherent
sheaf
\[
\mathcal{F}_{\mathrm{glob}}^{\mathrm{prism}}\;\in\;\mathrm{Coh}(\mathcal{X}^{\mathrm{Sel}}_{\bar\rho,\det,\mathcal{L}}).
\]
\end{theorem}

\begin{proof}
\emph{Step 1: Pairwise compatibility.} For each pair \(p\neq q\) of
finite primes, the pullback of
\(\mathcal{F}_{p}\boxtimes\mathcal{F}_{q}\) along the transition map
of \Cref{v4-arith:bcglue-k5:def:nygaard-transition} is a coherent
sheaf \(\mathcal{F}_{p,q}\) on \(\mathcal{X}^{\mathrm{EG},\{p,q\}}_{GL_{2}}\).
Prismatic flatness (\Cref{v4-arith:bcglue-k5:lem:prismatic-flatness})
ensures that the pullback preserves coherence.

\emph{Step 2: Triple compatibility.} For distinct \(p,q,r\), the two
constructions
\((\mathcal{F}_{p,q})\boxtimes\mathcal{F}_{r}\) and
\(\mathcal{F}_{p}\boxtimes(\mathcal{F}_{q,r})\), pulled back along
the respective transition maps into
\(\mathcal{X}^{\mathrm{EG},\{p,q,r\}}_{GL_{2}}\), agree by the
associativity of the Nygaard gluing data (Bhatt--Lurie 2022 \S8,
Cor.~8.2). Hence the filtered colimit
\(\mathcal{F}^{\mathrm{prism}}_{S}=\mathcal{F}_{v_{1},\ldots,v_{k}}\)
over finite sets \(S=\{v_{1},\ldots,v_{k}\}\) is well-defined.

\emph{Step 3: Archimedean inclusion.} At the archimedean place, the
local coherent Springer sheaf \(\mathcal{F}_{\infty}\) is the
Kottwitz--Langlands 1985 archimedean moduli sheaf on
\(\mathcal{X}^{\mathrm{EG},\infty}_{GL_{2}}\) (the moduli of rank-\(2\)
\(W_{\mathbb{R}}\)-representations). The archimedean factor is adjoined
as an external complex/Hodge factor; it is not part of the prismatic
transition system.
Including \(v=\infty\) in the filtered colimit is well-defined.

\emph{Step 4: Filtered colimit.} The filtered colimit
\[
\mathcal{F}_{\mathrm{glob}}^{\mathrm{prism}}\;:=\;\varinjlim_{S\text{ finite}}\mathcal{F}^{\mathrm{prism}}_{S}
\]
would exist in \(\mathrm{Coh}(\mathcal{X}^{\mathrm{Sel}}_{\bar\rho,\det,\mathcal{L}})\)
after proving the prismatic transition maps and the required
continuity of \(\mathrm{IndCoh}_{\mathrm{Nilp}}\) in the setting of
Gaitsgory--Rozenblyum's AMS DAG monographs. Thus (BCglue.glue) is
conditional, not closed.
\end{proof}

\subsection{Reduction to core compatibility}

\begin{theorem}[BC-diamond-glue reduces to core compatibility]
\label{v4-arith:bcglue-k5:thm:core-compatibility}
\ClaimStatusConjectured \emph{(conditional on arithmetic AG,
\Cref{v4-arith:thm:arithmetic-AG})}. After (BCglue.local) and
(BCglue.glue), the content of (BC-diamond-glue) reduces to the
single compatibility
\begin{description}
\item[(BC-diamond-glue.core)] \emph{Global Bernstein-centre
compatibility.} For every pair of distinct primes \(p\neq q\), the
diagonal pullback
\[
\Delta^{*}_{p,q}\mathcal{F}_{\mathrm{glob}}^{\mathrm{prism}}\;\in\;\mathrm{Coh}(\mathcal{X}^{\mathrm{EG},\{p,q\}}_{GL_{2}})
\]
of the glued sheaf onto the diagonal \((p,q)\)-stratum agrees with
\(\mathcal{F}_{p}\boxtimes\mathcal{F}_{q}\) under the action of the
tensor product of local Bernstein centres
\(\mathcal{Z}_{p}\otimes\mathcal{Z}_{q}\), and the induced
colimit
\(\varinjlim_{S}\mathcal{Z}_{S}=\mathcal{Z}_{\mathrm{glob}}\) agrees
with the global Bernstein centre acting on
\(\mathrm{D}^{\mathrm{sph\text{-}fin}}([GL_{2}])\) via the conjectural
arithmetic AG equivalence of \Cref{v4-arith:thm:arithmetic-AG}.
\end{description}
(BC-diamond-glue.core) is a statement of non-formal theorem: a
single cross-prime compatibility on the Selmer-derived global spectral stack.
\end{theorem}

\begin{proof}
By \Cref{v4-arith:bcglue-k5:thm:decomposition}, (BC-diamond-glue)
is the conjunction of three operations. (BCglue.local) is
conditional by \Cref{v4-arith:bcglue-k5:thm:local-BZCHN-plus-EG};
(BCglue.glue) is conjectural by
\Cref{v4-arith:bcglue-k5:thm:prismatic-gluing}. If both are supplied,
the remaining residual is (BCglue.core).

(BCglue.core) is the compatibility between two a priori independent
actions on \(\mathrm{IndCoh}_{\mathrm{Nilp}}(\mathcal{X}^{\mathrm{Sel}}_{\bar\rho,\det,\mathcal{L}})\):
the one from \(\mathcal{F}_{\mathrm{glob}}^{\mathrm{prism}}\) (via
Ben-Zvi--Chen--Helm--Nadler and prismatic gluing) and the one from the global Bernstein
centre (via the conjectural arithmetic AG). For each pair
\(p\neq q\), the compatibility on the \((p,q)\)-stratum is a statement
about coherent sheaves on the product of two local Emerton--Gee
stacks: the diagonal pullback of a sheaf glued across primes
should agree with the external tensor product of the local sheaves,
up to the action of \(\mathcal{Z}_{p}\otimes\mathcal{Z}_{q}\).

This pairwise compatibility is a finite pairwise check (standard
coherent-sheaf comparison along a flat gluing); taking the
filtered colimit over pairs to obtain the global compatibility is
non-formal theorem of standard filtered-colimit argument. The total
domain is non-formal theorem.
\end{proof}

\begin{remark}[(BC-diamond-glue.core) vs.\ unconditional closure]
\label{v4-arith:bcglue-k5:rem:core-vs-unconditional}
(BC-diamond-glue.core) is conditional on the arithmetic AG
equivalence \Cref{v4-arith:thm:arithmetic-AG}, which is itself a
global conjecture: Fargues--Scholze supplies the local form,
but the global arithmetic AG is not currently in the primary
literature. Hence (BC-diamond-glue.core) is \emph{conjectural
conditional on arithmetic AG}; it is not unconditional.

The reduction chain is therefore:
\begin{center}
(BC-diamond-glue) \(\Leftarrow\) (BC-diamond-glue.core) \(\Leftarrow\)
arithmetic AG.
\end{center}
The irreducible residual of (BC-diamond-glue) is
(BC-diamond-glue.core): a single cross-prime compatibility requiring
a non-formal theorem, conditional on arithmetic AG. The primary-source
conditions remaining are the local Emerton--Gee comparison, the
cross-prime prismatic transition maps, and arithmetic AG.
\end{remark}

\section{Consequence: \(\mathrm{L}_{\mathrm{ACD}}\) Final Residual Isolation}
\label{v4-arith:bcglue-k5:LACD-final}

\begin{corollary}[\(\mathrm{L}_{\mathrm{ACD}}\) tight residual list]
\label{v4-arith:bcglue-k5:cor:LACD-tight-list}
\ClaimStatusConjectured \emph{(conditional on prior F2.c, F1.a,
and on (BCglue.local) and (BCglue.glue) being proved)}. Under these
reductions, \(\mathrm{L}_{\mathrm{ACD}}\)
(\Cref{v4-arith:acd:thm:LACD-single-lemma})
has the following residual list:
\begin{enumerate}
\item (BCglue.local): \textbf{Conditional}; Ben-Zvi--Chen--Helm--Nadler supply the
local coherent Springer input, while the Emerton--Gee comparison and
\(\ell=p\) compatibility remain residual in
\Cref{v4-arith:bcglue-k5:thm:local-BZCHN-plus-EG}.
\item (BCglue.glue): \textbf{Conjectural}; no primary-source
cross-prime prismatic transition theorem currently supplies
\Cref{v4-arith:bcglue-k5:thm:prismatic-gluing}.
\item (BC-diamond-glue.core): \textbf{Conjecture conditional on
arithmetic AG} (\Cref{v4-arith:thm:arithmetic-AG}), isolated in
\Cref{v4-arith:bcglue-k5:thm:core-compatibility}.
\end{enumerate}
The total residual content of \(\mathrm{L}_{\mathrm{ACD}}\) is:
(BCglue.local) \(+\) (BCglue.glue) \(+\) (BC-diamond-glue.core), with
the core item conditional on arithmetic AG. Required theorem past
arithmetic AG: non-formal theorem.
\end{corollary}

\begin{proof}
Assemble \Cref{v4-arith:bcglue-k5:thm:decomposition},
\Cref{v4-arith:bcglue-k5:thm:local-BZCHN-plus-EG},
\Cref{v4-arith:bcglue-k5:thm:prismatic-gluing}, and
\Cref{v4-arith:bcglue-k5:thm:core-compatibility}.
(BCglue.local) is conditional; (BCglue.glue) is conjectural;
(BC-diamond-glue.core) is the residual after those inputs are supplied.
\end{proof}

\begin{remark}[domain decomposition of (BC-diamond-glue)]
\label{v4-arith:bcglue-k5:rem:domain-delta}
\Cref{v4-arith:acd:rem:bc-diamond-glue-domain} estimated
(BC-diamond-glue) at non-formal theorem past Ben-Zvi--Chen--Helm--Nadler.
The decomposition of \Cref{v4-arith:bcglue-k5:thm:decomposition}
partitions that domain into non-formal theorem for (BCglue.local), non-formal theorem
for (BCglue.glue), and non-formal theorem for (BC-diamond-glue.core),
conditional on arithmetic AG. None of these three pieces is closed
by the preceding inputs.
\end{remark}

\section{K5: Godement--Jacquet Absorption}
\label{v4-arith:bcglue-k5:K5}

\subsection{The Local Adjoint \(L\)-Factor}

\begin{definition}[local adjoint \(L\)-factor \(L_{p}(s,\mathrm{Ad})\)]
\label{v4-arith:bcglue-k5:def:adjoint-L}
For an unramified principal-series representation
\(\pi_{p}=\mathrm{Ind}_{B}^{GL_{2}}(\chi_{1}\otimes\chi_{2})\) at
\(GL_{2}(\mathbb{Q}_{p})\) with Satake parameters
\((\alpha_{p},\beta_{p})=(\chi_{1}(p),\chi_{2}(p))\) on the unit
circle, the local adjoint \(L\)-factor is
\begin{equation}
\label{v4-arith:bcglue-k5:eq:adjoint-L}
L_{p}(s,\mathrm{Ad}(\pi_{p}))
\;:=\;\bigl(1-\alpha_{p}\beta_{p}^{-1}\,p^{-s}\bigr)^{-1}
\cdot(1-p^{-s})^{-1}
\cdot\bigl(1-\alpha_{p}^{-1}\beta_{p}\,p^{-s}\bigr)^{-1}.
\end{equation}
Reference: Jacquet 1979 Corvallis ``Principal \(L\)-functions'' \S2,
or Bump 1997 \emph{Automorphic Forms and Representations} \S3.2.
\end{definition}

\begin{lemma}[the regular domain of the local adjoint scalar]
\label{v4-arith:bcglue-k5:lem:Lp-1-well-defined}
For specified nonzero unramified parameters \(\alpha_p,\beta_p\),
put \(r=\alpha_p\beta_p^{-1}\). If
\(p^{-1}<|r|<p\), then the displayed scalar
\(L_p(1,\mathrm{Ad})\) is finite and nonzero. It is positive real
when either \(|r|=1\), or \(r\) is real with \(p^{-1}<r<p\).
For general \(r\ne0\), its only poles as a function of \(r\) occur
at \(r=p\) and \(r=p^{-1}\).
\end{lemma}

\begin{proof}
The three denominator factors are
\[
 1-r/p,\qquad 1-p^{-1},\qquad 1-r^{-1}/p.
\]
In the stated annulus, both \(|r|/p\) and \(|r|^{-1}/p\) are
strictly smaller than one, so no factor vanishes. More explicitly,
\[
 |L_p(1,\mathrm{Ad})|
 \leq\frac{1}{(1-|r|/p)(1-p^{-1})(1-|r|^{-1}/p)}.
\]
This follows by applying \(|1-z|\geq1-|z|\) to the first and
third factors. If \(|r|=1\), then \(r^{-1}=\overline r\), so their
product is \(|1-r/p|^2>0\). If \(p^{-1}<r<p\) is real, both
factors are positive separately. The middle factor is always positive.
This proves the positivity statements. For arbitrary \(r\ne0\),
a zero of the first or third factor gives exactly the two stated
values; each gives a pole since the other factors are nonzero.
\end{proof}

An application to an automorphic representation must first supply its
parameters in this range. Tempered parameters satisfy \(|r|=1\).
More generally, bounds \(|\alpha_p|,|\beta_p|\leq p^\theta\)
with \(|\alpha_p\beta_p|=1\) and \(\theta<1/2\) imply
\(p^{-2\theta}\leq|r|\leq p^{2\theta}\), which is inside the
regular annulus. This implication is a scalar estimate, not a
construction assigning an automorphic representation to a native
Hochschild class. Ramified parameters require their actual local
factor. The pairing comparison below therefore starts with a
specified nonzero scalar and states positivity separately.

\subsection{The normalization and its comparison maps}

\begin{definition}[rescaling a specified Koszul pairing]
\label{v4-arith:bcglue-k5:def:koszul-dual-norm}
Let \(V_p,V_p^!\) be the native carriers, \(B_p:V_p\times V_p^!\to
\mathbb C\) their bilinear pairing, and \(\lambda_p\in\mathbb C^\times\)
a specified scalar. Define
\begin{equation}
             B_p^{\mathrm{Ad}}=\lambda_p^{-1}B_p.
             \label{v4-arith:bcglue-k5:eq:koszul-dual-norm}
\end{equation}
The adjoint-factor convention takes \(\lambda_p=L_p(1,\mathrm{Ad})\)
where that value is finite and nonzero. To divide a positive Hermitian
form by the same scalar and retain positivity, one must additionally
have \(\lambda_p>0\).
\end{definition}

\begin{conjecture}[native local comparison with the adjoint normalization]
\label{v4-arith:bcglue-k5:conj:local-comparison-datum}
On the specified cuspidal spherical domain, the original native pairing
and Whittaker form admit the following realization with their declared
Haar normalizations. There are linear maps
\(f_p:V_p\to\mathrm{HS}(\ell^2(\mathbb N_0))\),
\(g_p:V_p^!\to\mathrm{HS}(\ell^2(\mathbb N_0))\) and a
conjugate-linear map \(c_p:V_p\to V_p^!\), preserving the specified
arithmetic operations on their domains, such that
\[
 B_p(v,w^!)=\lambda_p B_q(f_pv,g_pw^!),\qquad
 h_p^{\mathrm{Whit}}(v,w)=\lambda_p h_q(f_pv,f_pw),\qquad
 g_pc_p=J_qf_p,
 \quad q=p^{-1}.
\]
Every compact-group volume factor must be fixed in these maps or forms.
For example the map \(\mathbf1_{GL_2(\mathbb Z_p)}\mapsto
\sqrt{(1-p^{-1})(1-p^{-2})}\,\Omega_p\) has precisely that Haar mass;
the unit state itself still has mass one. The conjecture concerns
actual maps on native carriers, not merely the equality of one scalar
spherical value.
\end{conjecture}

\begin{theorem}[local transfer under the native comparison]
\label{v4-arith:bcglue-k5:thm:K5-absorption-canonical}
Under \Cref{v4-arith:bcglue-k5:conj:local-comparison-datum}, define
\(h_p^{\mathrm{BC}}(v,w)=h_q(f_pv,f_pw)\). Then
\begin{equation}
 h_p^{\mathrm{BC}}(v,w)=\lambda_p^{-1}h_p^{\mathrm{Whit}}(v,w)
                       =B_p^{\mathrm{Ad}}(v,c_pw).
 \label{v4-arith:bcglue-k5:eq:K5-match}
\end{equation}
If \(f_p\) is injective, this form is positive definite. These
identities hold on the common specified domain of the maps.
\end{theorem}
\begin{proof}
The first equality is the assumed Whittaker realization. For the second,
\[
 B_p^{\mathrm{Ad}}(v,c_pw)=B_q(f_pv,g_pc_pw)
  =B_q(f_pv,J_qf_pw)=h_q(f_pv,f_pw),
\]
where the final equality is the explicit Hilbert--Schmidt calculation
of \Cref{v4-arith:thermal:bilinear-pairing}. Its diagonal value is
\(\|f_pv\|_2^2\), which is positive for \(v\ne0\) when \(f_p\)
is injective. Both \(h_p^{\mathrm{Whit}}\) and \(h_p^{\mathrm{BC}}\)
are Hermitian; inserting \(J_q\) into the second slot of either would
instead make the expression bilinear. Modular conjugation occurs in
the bilinear form \(B_q\).
\end{proof}

\begin{remark}[what a spherical scalar can establish]
\label{v4-arith:bcglue-k5:rem:K5-domain}
An equality on one vacuum line does not identify two forms. On
\(\mathbb C^2\), the positive forms with matrices
\(\mathrm{diag}(1,1)\) and \(\mathrm{diag}(1,2)\) agree on the first
coordinate line and disagree on the second. The same test applies to
an attempted deduction from a spherical integral to all native vectors.
A local functional equation is a statement about specified integrals
and parameters; it does not supply the maps \(f_p,g_p,c_p\).
\end{remark}

\begin{remark}[scalar factors and measures]
\label{v4-arith:bcglue-k5:eq:GJ-funeq}
The volume \(1-p^{-1}\) and the adjoint factor are separate scalars.
For instance the algebraic identity
\(L_p(s,\pi\times\widetilde\pi)=\zeta_p(s)L_p(s,\mathrm{Ad})\)
would give
\(L_p(1,\pi\times\widetilde\pi)\zeta_p(1)^{-1}
 =L_p(1,\mathrm{Ad})\), wherever these factors are finite. It does
not leave another factor \(1-p^{-1}\). A further volume can enter
only through an independently specified measure or vector scaling.
\end{remark}

\subsection{Adelic normalization}

\begin{theorem}[conditional adelic rescaling]
\label{v4-arith:bcglue-k5:thm:K5-adelic-consistent}
Assume the native local comparisons, their compatible finite-prime
assembly, and a common adelic regularization of their forms and scalar
factors. If this construction gives a finite nonzero scalar
\(\lambda\), then rescaling both assembled forms gives
\[
       h^{\mathrm{Whit,reduced}}=\lambda^{-1}h^{\mathrm{Whit}},
       \qquad B^{\mathrm{Ad}}=\lambda^{-1}B,
       \qquad h^{\mathrm{Whit,reduced}}(v,w)=B^{\mathrm{Ad}}(v,cw).
\]
If the unreduced Hermitian form is positive and \(\lambda>0\), the
reduced form is positive on the same injective native realization.
The choice \(\lambda=L(1,\mathrm{Ad})\) requires identification of
that continued value with the scalar furnished by the stated assembly.
\end{theorem}
\begin{proof}
For a finite prime set the local identities multiply, including every
scalar and volume factor. The assumed common assembly preserves these
identities. Division by the nonzero assembled scalar gives the stated
equality, and division by a positive real scalar preserves positivity.
Analytic continuation to \(s=1\) is not a proof of absolute Euler-product
convergence at that point and is not a construction of the native
assembly. Neither is inferred in this argument.
\end{proof}

\begin{remark}[local scalars and global assembly]
\label{v4-arith:bcglue-k5:rem:K5-delta}
Once a nonzero scalar is specified, dividing a bilinear pairing by
that scalar is a definition. Equality with a second pairing requires
the local maps of
\Cref{v4-arith:bcglue-k5:conj:local-comparison-datum}.
A global division additionally requires the assembled scalar in
\Cref{v4-arith:bcglue-k5:thm:K5-adelic-consistent}.
Neither condition is removed by nonvanishing of an analytic
\(L\)-value.
\end{remark}

\section{Consequence: K-P Compat, Conditional Domain}
\label{v4-arith:bcglue-k5:K-P-unconditional}

\begin{corollary}[conditional global pairing consequence]
\label{v4-arith:bcglue-k5:cor:K-P-unconditional}
Assume the finite native comparisons, the archimedean pairing maps of
\Cref{v4-arith:kp:prop:Pet-is-Segal-infty}, and compatible adelic
assembly. Then the native Koszul--Petersson identity holds on their
common domain. If an injective automorphic \(L^2\) realization is also
supplied, its pulled-back norm is positive definite.
\end{corollary}
\begin{proof}
Apply the finite and archimedean transfer identities, then the assumed
common assembly. Positivity follows from the injective realization into
a positive Hilbert space. The local normalization does not construct
any of these native maps.
\end{proof}

\begin{remark}[what the finite comparison supplies]
\label{v4-arith:bcglue-k5:rem:P3-impact}
The finite K5 comparison supplies one part of the pairing problem.
Archimedean conjugation and native injectivity are separate conditions;
positive local states do not imply an injective global arithmetic map.
The variance of the native duality must be preserved at every place.
\end{remark}

\begin{remark}[spectral and positivity requirements]
\label{v4-arith:bcglue-k5:rem:RH-impact}
A sufficient arithmetic implication to the Riemann hypothesis retains
its native duality, pairing, zero-divisor and positivity hypotheses.
The local scalar rescaling proves no Hilbert--P\'olya realization and
removes no independent arithmetic trace comparison.
\end{remark}

\section{Residual Conditions}
\label{v4-arith:bcglue-k5:residual}

The open residuals split as follows.

\begin{description}
\item[(BCglue.local).] Conditional;
\Cref{v4-arith:bcglue-k5:thm:local-BZCHN-plus-EG}.
Residual condition: non-formal theorem.
\item[(BCglue.glue).] Conjectural (prismatic flatness);
\Cref{v4-arith:bcglue-k5:thm:prismatic-gluing}.
Residual condition: non-formal theorem.
\item[(BC-diamond-glue.core).] Residual conjecture conditional on
arithmetic AG (\Cref{v4-arith:thm:arithmetic-AG}). Domain:
non-formal theorem past arithmetic AG.
\item[Arithmetic AG (global).] Open conjecture. Domain: substantial
(primary-literature-length); this is
the cumulative domain of extending Fargues--Scholze local
geometrization to a global arithmetic statement, estimated
non-formal theorem past Fargues--Scholze.
\item[K5 (GJ absorption convention).] Conditional on
\Cref{v4-arith:bcglue-k5:conj:local-comparison-datum} on the cuspidal
sph-fin locus by
\Cref{v4-arith:bcglue-k5:thm:K5-absorption-canonical} and
\Cref{v4-arith:bcglue-k5:thm:K5-adelic-consistent}. Remaining condition:
the local maps and compatible adelic assembly. Eisenstein handling via the
sph-fin cut of \Cref{v4-arith:sw:thm:KP-split}(K7).
\item[K-P compat, P3.] Domain-restricted to the cuspidal sph-fin
locus by \Cref{v4-arith:bcglue-k5:cor:K-P-unconditional}. Domain
conditional on the local comparison datum. The full-sector extension
has the KMS\(_{1}\)/GNS half-density factor separated from the
ramified Rankin--Selberg and Eisenstein conditions.
\item[Sufficient RH path.] Four-input conditional conjunction
(\Cref{v4-arith:w3-comparison-B:thm:four-input-branch}): arith.~Positselski
(five open sub-items), Koszul--Pet compat (domain-restricted),
H-P\(^{\mathrm{Ar}}\), VB\(^{*}\) main direction.
Arith.~Positselski: \(\sim 30\)--\(50\) hypothesis.
VB\(^{*}\) is the classical positive-trace open problem.
\end{description}

\begin{remark}[residual form]
\label{v4-arith:bcglue-k5:rem:residual-summary}
The residual form is:
\begin{enumerate}
\item (BCglue.local): non-formal theorem, form \emph{conditional}; the
integration of Ben-Zvi--Chen--Helm--Nadler + Emerton--Gee + Zhu +
Fargues--Scholze into a single canonical local coherent Springer sheaf,
including \(\ell=p\) compatibility and \(\ell\)-independence, remains a
primary-source writing condition; see
\Cref{v4-arith:bcglue-k5:thm:local-BZCHN-plus-EG}.
\item (BCglue.glue): non-formal theorem, form \emph{conjectural} conditional
on the prismatic flatness of
\Cref{v4-arith:bcglue-k5:lem:prismatic-flatness}; the cross-prime
Nygaard-convention-aware transition theorem is not in the primary
literature.
\item K5 (GJ absorption): reduced on the cuspidal sph-fin locus to
the local maps and compatible adelic assembly under the scalar convention
(\Cref{v4-arith:bcglue-k5:def:koszul-dual-norm}); Eisenstein
handling via the sph-fin cut.
\item P3 on the native cuspidal core requires the injective
automorphic realization, archimedean and finite pairing maps and
compatible adelic assembly of the conditional corollary.
\item (BC-diamond-glue): a single cross-prime compatibility remains,
non-formal theorem past arithmetic AG.
\item Arithmetic AG, arithmetic Positselski, and VB\(^{*}\) are not
logical consequences of K5.
\end{enumerate}
\end{remark}

\section{Residual Boundary}
\label{v4-arith:bcglue-k5:summary}

\begin{enumerate}
\item (BC-diamond-glue) decomposes into three sub-operations:
(BCglue.local), (BCglue.glue), (BC-diamond-glue.core)
(\Cref{v4-arith:bcglue-k5:thm:decomposition}). The decomposition
is logically independent.
\item (BCglue.local) is conditional: Ben-Zvi--Chen--Helm--Nadler +
Emerton--Gee + Zhu + Fargues--Scholze supply the inputs, but a single
canonical primary-source transition theorem covering \(\ell = p\) and
\(\ell\)-independence has not yet been written
(\Cref{v4-arith:bcglue-k5:thm:local-BZCHN-plus-EG}).
\item (BCglue.glue) is conjectural conditional on prismatic flatness
of the Nygaard-convention-aware transition maps
(Bhatt--Scholze, Bhatt--Lurie; see
\Cref{v4-arith:bcglue-k5:lem:prismatic-flatness}) and filtered
colimits (Gaitsgory--Rozenblyum)
(\Cref{v4-arith:bcglue-k5:thm:prismatic-gluing}).
\item The irreducible residual of (BC-diamond-glue) is
(BC-diamond-glue.core): a non-formal theorem cross-prime compatibility,
conditional on arithmetic AG
(\Cref{v4-arith:bcglue-k5:thm:core-compatibility}).
\item \(\mathrm{L}_{\mathrm{ACD}}\) tight residual list
(\Cref{v4-arith:bcglue-k5:cor:LACD-tight-list}): (BCglue.local)
conditional; (BCglue.glue) conjectural; (BC-diamond-glue.core)
conjectural conditional on arithmetic AG. The three subconditions
together form the residual content of \(\mathrm{L}_{\mathrm{ACD}}\)
past F2.c and the subconditions of
Ch.~\ref{v4-arith:lacd-sub}.
\item K5 requires the local maps of
\Cref{v4-arith:bcglue-k5:conj:local-comparison-datum} and the
compatible adelic assembly of
\Cref{v4-arith:bcglue-k5:thm:K5-adelic-consistent}.
At a finite set of primes the scalar rescaling is the product of
\(L_p(1,\mathrm{Ad})^{-1}\). Its extension to all primes and
identification with \(L(1,\mathrm{Ad})^{-1}\) require the
stated convergence or regularization. Eisenstein factors require
separate domains and measures.
\item K-P compat and P3 are conditional domain-restricted statements
on the cuspidal sph-fin locus with the explicit normalisation convention
(\Cref{v4-arith:bcglue-k5:cor:K-P-unconditional});
the full-sector extension separates the KMS\(_{1}\)/GNS half-density
from the ramified Rankin--Selberg and Eisenstein conditions.
\item Residual condition: (BC-diamond-glue) has one irreducible
cross-prime compatibility past arithmetic AG; the primary open
residuals remain arithmetic AG, arith.~Positselski (five open
sub-items), H-P\(^{\mathrm{Ar}}\), and VB\(^{*}\) (the
four-input conditional conjunction of
\Cref{v4-arith:w3-comparison-B:thm:four-input-branch}).
\end{enumerate}

(BC-diamond-glue) is isolated as a single cross-prime compatibility
conditional on arithmetic AG, and K5 is reduced on the cuspidal
sph-fin locus to native local maps and compatible assembly. The K-P compatibility
and P3 are conditional domain-restricted statements on the cuspidal
sph-fin locus. The sufficient RH path is unaffected.
